\chapter{The derived category of a dg-category}
\label{ch:derived}

This chapter is Keller \S 3.2 up to the definition, together with the comparison theorem
that is the finish line of the first phase. Keller notes that from the definition alone it is
not clear that $\Dcat(\cA)$ has small hom-sets or is additive; both facts are established
classically through resolutions (Proposition~\ref{prop:resolutions}), which is why that
proposition sits on the critical path even though model structures do not.

\begin{definition}[Quasi-isomorphisms of dg-modules]
\label{def:dg-qis}
% planned: DgCat.DGModule.quasiIso
\uses{def:dg-module-cat, def:dg-module-homology}
$\qis$ is the class of morphisms of $\Ch(\cA)$ inducing isomorphisms on all
$H^n(M(X))$. It is a multiplicative class containing the isomorphisms and closed under
composition and retracts.
\end{definition}

\begin{definition}[The derived category]
\label{def:dg-derived-category}
% planned: DgCat.DerivedCategory, DgCat.DerivedCategory.Q
\uses{def:dg-qis, def:localization}
$\Dcat(\cA)$ is the localization of $\Ch(\cA)$ at $\qis$, with projection functor
$Q : \Ch(\cA) \to \Dcat(\cA)$, constructed through mathlib's localization framework.
\end{definition}

\begin{lemma}[Localizing the homotopy category instead]
\label{lem:derived-via-homotopy}
% planned: DgCat.DerivedCategory.ofHomotopyCategory
\uses{def:dg-derived-category, def:dg-module-homotopy-cat}
Null-homotopic morphisms become zero in $\Dcat(\cA)$, so $Q$ factors through
$\Hcat(\cA)$, and $\Dcat(\cA)$ is also the localization of $\Hcat(\cA)$ at the image of
$\qis$.
\end{lemma}
\begin{proof}
A null-homotopic map factors through a contractible complex, which is quasi-isomorphic to
zero. Compare mathlib's proof of the same statement for complexes in an abelian category.
\end{proof}

\begin{definition}[Cofibrant and fibrant dg-modules]
\label{def:cofibrant-fibrant}
% planned: DgCat.DGModule.IsCofibrant, DgCat.DGModule.IsFibrant
\uses{def:dg-module-cat, def:dg-qis}
A dg-module $P$ is cofibrant if every morphism $P \to M$ factors through every surjective
quasi-isomorphism $L \to M$; $I$ is fibrant if every morphism $L \to I$ extends along every
injective quasi-isomorphism $L \to M$. Representables $X^{\wedge}$ are cofibrant.
\end{definition}

\begin{proposition}[Resolutions]
\label{prop:resolutions}
% planned: DgCat.DGModule.cofibrantResolution, DgCat.DGModule.fibrantResolution
\uses{def:cofibrant-fibrant, def:dg-module-homotopy-cat, lem:dg-homotopy-yoneda}
For every dg-module $M$ there is a quasi-isomorphism $\mathbf{p}M \to M$ with
$\mathbf{p}M$ cofibrant and a quasi-isomorphism $M \to \mathbf{i}M$ with $\mathbf{i}M$
fibrant. The projection $\Hcat(\cA) \to \Dcat(\cA)$ has a fully faithful left adjoint
$M \mapsto \mathbf{p}M$ and a fully faithful right adjoint $M \mapsto \mathbf{i}M$
(Keller, Proposition 3.1).
\end{proposition}
\begin{proof}
\uses{lem:derived-via-homotopy}
Keller's explicit construction (after Keller 1994, \S 3), by transfinite induction on
semi-free filtrations; the alternative through the projective model structure
(Chapter~\ref{ch:beyond}) is deferred. This is the largest single proof in the first phase.
\end{proof}

\begin{corollary}[Hom-sets in the derived category]
\label{cor:derived-homs}
% planned: DgCat.DerivedCategory.homEquiv
\uses{prop:resolutions, def:dg-derived-category}
$\Dcat(\cA)(L,M) \cong \Hcat(\cA)(\mathbf{p}L, M) \cong \Hcat(\cA)(L, \mathbf{i}M)$.
In particular $\Dcat(\cA)$ has small hom-sets and is additive, and for an object $X$,
$\Dcat(\cA)(X^{\wedge}, M[n]) \cong H^n M(X)$ (Keller (3)).
\end{corollary}
\begin{proof}
\uses{lem:dg-homotopy-yoneda}
Fully faithful adjoints; then Lemma~\ref{lem:dg-homotopy-yoneda} since $X^{\wedge}$ is
cofibrant.
\end{proof}

\begin{theorem}[Comparison with the derived category of an algebra]
\label{thm:ring-comparison}
% planned: DgCat.DerivedCategory.ringComparison
\uses{def:dg-derived-category, prop:dg-module-ring, def:mathlib-derived-category, cor:derived-homs}
If $\cA = \cA_B$ is the one-object dg-category of a $k$-algebra $B$, then the isomorphism
$\Ch(\cA) \cong \Ch(B)$ of Proposition~\ref{prop:dg-module-ring} identifies
quasi-isomorphisms on both sides and induces an equivalence of categories
$\Dcat(\cA) \simeq \Dcat(\Mod B)$ with mathlib's derived category, compatible with the
projection functors.
\end{theorem}
\begin{proof}
\uses{prop:dg-module-ring, def:localization}
Both sides are localizations of isomorphic categories at corresponding classes of
morphisms; the universal property of localization gives the equivalence. The content is in
checking that the two notions of quasi-isomorphism match under
Proposition~\ref{prop:dg-module-ring}, which is where a wrong sign or a wrong shift
convention would surface.
\end{proof}

\begin{remark}
This theorem is the acceptance test for the whole design. If the enriched-category
definition of a dg-category, the transport-of-enrichment construction of $\Zz$, and the
localization-based derived category did not reproduce mathlib's $\Dcat(\Mod B)$, the
infrastructure would be of no use downstream.
\end{remark}
